\section{Discussion}

\paragraph{An effect on a receiver is not a signal.}
The operational definition used by \citet{singh2025active}---an \eod\ that
conspecifics detect and that changes their behaviour---is satisfied in our model
at every metabolic cost we tested, including zero. It is satisfied because the
knollenorgan channel really does carry information and receivers really do
respond to it. Yet the same channel has no detectable consequence for anyone's
payoff: replaying, scrambling or fabricating it moves neither party's return by
an amount we can resolve. We are careful about what that licenses. Our
equivalence tests do not clear the \SI{5}{\percent} margin we set in advance,
because the confidence intervals are wider than the margin; the defensible
statement is \emph{no effect detectable at this sample size}, not demonstrated
equivalence. What makes the negative interpretable is not the width of those
intervals but the assay validation of \S\ref{sec:assay}: the same tests, applied
to a dyad that is communicating by construction, return effects one to two orders
of magnitude larger than anything we see here.

So the strongest defensible form of the claim is this. A system can satisfy the
standard positive-signalling and positive-listening diagnostics without any
detectable payoff relevance and without evidence that production was shaped by
being heard.

Establishing the production-side half of the argument turned out to need more
care than we first gave it. Disabling knollenorgans removes \emph{reception}, so
a hearing/deaf contrast asks whether a fish behaves differently when it can hear
others---not whether it behaves differently when others can hear \emph{it}. Those
are different questions, and the first is contaminated by free-riding: a fish
that can hear its neighbours' probes cuts its own costly sampling, which is what
we in fact observe. The control that isolates being listened to has to hold
reception fixed and cut audibility, which is what our yoked condition does.

\paragraph{The muting confound.}
Silencing a fish is the natural intervention and it is the one available in real
experiments, but it removes three things simultaneously. Our decomposition shows
that the three shares are not merely different in size---they act on different
individuals and in different directions. Losing reafference hurts the muted fish;
losing the illumination it provided hurts its neighbours; losing its detectability
changes how the group is spaced. A single number attributed to ``communication''
averages over all three. The same caution applies to real ablation work: an
\eod-silenced \emph{Gnathonemus} is not a fish with its telephone taken away, it is
a fish that has been blinded and has stopped lighting the room for everyone else.

\paragraph{Cost does not manufacture semantics.}
It is natural to expect that pricing a pulse would turn a probe schedule into an
informative one: a budget has to be spent where it pays, and that allocation
should make the train diagnostic of the emitter's state. That is not what
happens. Across the whole cost range the pulse train becomes rarer and
\emph{less} informative about every target we decoded, and it stays less
informative after normalising per pulse. The learned response to a price is a
lower baseline rate, not a smarter schedule. Whatever costly signalling does in
the classical models
\citep{zahavi1975mate,grafen1990biological,szamado2011cost}, it presupposes a
signal whose honesty is at stake; it does not by itself bring one into being.

What we should not conclude is that a binary pulse leaves no room for a code. A
binary event repeated in continuous time supports a large space of temporal
codes---inter-pulse intervals, bursts, pauses, event-locked patterns---and
mormyrid communication is known to use pulse timing and waveform identity
\citep{baker2013multiplexed,carlson2004stereotyped,hopkins1988neuroethology}. Our
agents \emph{failed to exploit} the temporal code available to them under the
pressures we applied. That is a claim about what learning found, not about what
the channel permits.

\paragraph{Silence tracks the regime, not the moment.}
Cessation of discharge is documented in mormyrids as both predator avoidance and
submissive display \citep{stoddard1999predation,moller1989electric}, and we
expected a momentary version: quiet when the predator is near. It does not appear.
Discharge probability is essentially unchanged by a predator within
\SI{25}{\centi\metre} (\SilenceDangerDiff) and danger is near-undecodable from the
public channel (\DangerAUC\ AUC). Agents lower their \emph{baseline} rate wherever
predators exist, and mortality falls with it, so crypsis works without timing. The
likely reason is sensory asymmetry: our predator hears discharges at
\SI{60}{\centi\metre} but is passively detectable only at \SI{25}{\centi\metre}.

\paragraph{Signalling as a public good and a public hazard.}
Two externalities operate in opposite directions, neither in the reward function:
a discharge illuminates neighbours' worlds, and it also jams their passive sense
and draws an eavesdropping predator toward the group. Cooperative harvesting
sustains a higher rate than competition at matched cost, consistent with the
public-good component being internalised when the harvest is shared. We do not
state that as a mechanism: our interventions found no detectable payoff
consequence of the public channel, so we cannot also claim neighbours benefit
enough from a pulse to explain the rate difference, which may instead reflect
altered encounter distributions or an effect below our resolution.

\paragraph{What the discharge lacks is not a channel.}
The decoupled control was meant to test the hypothesis that dual use is what
blocks signalling: give the fish a variable it can move without degrading its own
sensing, and a social function should appear. It does not. Agents mark half their
discharges with the subtype---the rate an unmodulated bit would give---and
deleting it costs nobody anything in any of the four ecological settings.

Setting that beside the referential assay says what the missing ingredient is.
The assay differs from the foraging world in three ways, and only one of them is
the channel. It also gives the sender information the receiver cannot get for
itself, and it pays both of them when the receiver acts on it. In the foraging
world a fish can find food with its own probe, so a neighbour's private state is
worth very little; there is a channel available and nothing to say over it.
Channel factorisation is plausibly necessary for a signal to become distinct from
a probe, but it is clearly not sufficient. That is a testable prediction for the
biology: electric communication should be elaborated in exactly those contexts
where one fish reliably knows something a conspecific cannot discover
first-hand---dominance status, reproductive state, individual identity---and not
in contexts where both parties can simply sense the world.

\paragraph{Relation to emergent-communication work.}
Most emergent-communication studies equip agents with a dedicated, costless,
otherwise-useless channel \citep{foerster2016learning,mordatch2018emergence,lazaridou2017multi}.
Under those conditions any use of the channel is communicative by construction,
and the interesting question is what code emerges. Electric fish invert the setup:
the channel is a physical act with a large private benefit, and the interesting
question is whether \emph{any} of its use is communicative. We think this is the
more common situation in biology---most signals are thought to be ritualised from
behaviours that originally had other functions---and it is a setting in which the
standard metrics need the production-side test that SSI supplies. Positive
signalling and positive listening are both necessary; neither, nor both together,
is sufficient.

\section{Limitations}

Six bound the claims above and are set out in full in
Appendix~\ref{app:limitations}. The three that matter most: our populations adapt
by policy-gradient learning within a training run, not by an evolutionary outer
loop, so the Sender Shaping Index is a within-lifetime analogue of
\citeauthor{scottphillips2008defining}'s criterion rather than the criterion
itself, and we write \emph{emerge under learning} throughout. Our knollenorgan
channel is all-or-none, as upstream, so graded amplitude signalling is
structurally impossible in the coupled condition, and the decoupled condition
lifts only the crudest version of that restriction --- one bit, not a waveform.
And the equivalence tests do not reach the margin we set in advance; what makes
the nulls interpretable is the assay validation of \S\ref{sec:assay}, not the
width of those intervals.
